% Reconstructed text from V5 PDF page 41; non-native source
\par\smallskip\noindent\rule{\textwidth}{0.4pt}\par{\sffamily\footnotesize Original V5 page 41 -- recovered text}\par\begin{flushleft}\footnotesize\ttfamily
\par\smallskip
\par\smallskip
\noindent Diagnostics and Serving Infrastructure for MAS 41\par
\par\smallskip
\par\smallskip
\noindent Table 9: Retention ablation at 10\% budget on 220 TraceElephant traces (5 task suites, 3 MAS frame-\par
\noindent works). S2 shows mean ± SD across 3 seeds; S1, S3, S4 are deterministic. Higher is better.\par
\par\smallskip
\noindent  Strategy Causal Last-Error\par
\noindent  S1 Full retention 0.459 0.432\par
\noindent  S2 Random 0.524 ± 0.023 0.523 ± 0.020\par
\noindent  S3 Recency 0.391 0.391\par
\noindent  S4 Importance 0.532 0.500\par
\par\smallskip
\noindent Table 10: Causal-graph attribution accuracy across all retention budgets (220 TraceElephant traces). S1\par
\noindent is constant by definition. S2 shows seed mean. The inverted slope (accuracy decreasing toward 100\%) is\par
\noindent visible for S2 and S4.\par
\par\smallskip
\noindent  Budget S1 Full S2 Random S3 Recency S4 Importance\par
\noindent  10\% 0.459 0.524 0.391 0.532\par
\noindent  20\% 0.459 0.479 0.423 0.555\par
\noindent  30\% 0.459 0.479 0.423 0.546\par
\noindent  40\% 0.459 0.479 0.445 0.514\par
\noindent  50\% 0.459 0.461 0.450 0.496\par
\noindent  60\% 0.459 0.450 0.459 0.491\par
\noindent  70\% 0.459 0.432 0.459 0.482\par
\noindent  80\% 0.459 0.451 0.464 0.473\par
\noindent  90\% 0.459 0.457 0.459 0.464\par
\noindent  100\% 0.459 0.459 0.459 0.459\par
\par\smallskip
\par\smallskip
\noindent agent if all errors appear recovered); and (b) a last-error heuristic baseline that attributes to the\par
\noindent agent of the last error event in the retained window.\par
\noindent  The results show three patterns:\par
\noindent  1. Subsampling paradoxically improves heuristic accuracy. At 10\% retention budget,\par
\noindent  both random sampling (S2, 0.524) and importance-weighted sampling (S4, 0.532) exceed\par
\noindent  full-retention accuracy (S1, 0.459) under the causal-graph heuristic. The full budget curve\par
\noindent  (Table 10) shows accuracy decreasing as budget approaches 100\%, revealing a limitation\par
\noindent  of heuristic attributors rather than a universal benefit of less data.\par
\noindent  2. The mechanism: noise removal exposes the true root cause. The causal-graph heuris-\par
\noindent  tic marks an error as “recovered” if the same agent later produces a normal event. On\par
\noindent  full traces, the mistake agent typically continues producing normal events after its error,\par
\noindent  causing the heuristic to skip the true root cause. Subsampling strips these post-error nor-\par
\noindent  mal events, making the decisive error visible. Importance-weighted sampling amplifies\par
\noindent  this effect by systematically concentrating error events.\par
\noindent  3. Importance scoring adds only marginal value over random at low budgets. At 10\%\par
\noindent  budget, the S4–S2 gap is +0.008, within noise. At intermediate budgets (20–40\%), S4\par
\noindent  shows a more consistent advantage (Table 10). This suggests the primary benefit at ex-\par
\noindent  treme subsampling derives from noise removal rather than from the specific importance\par
\noindent  scores. Recency-based retention (S3, 0.391 at 10\%) is worst at low budgets because fail-\par
\noindent  ures whose root cause occurs early in the trace are systematically missed.\par
\noindent  Table 9 reports the verified results at 10\% retention budget. Table 10 shows the full causal-\par
\noindent method curve across all 10 budget levels.\par
\noindent  At full retention, the causal method outperforms last-error by 2.7 pp (0.459 vs. 0.432). This\par
\end{flushleft}
